\guidechapter{Why improving one part may change nothing}{The final saving is limited by the weakest margin.}

Suppose three operations have normalized costs $p^{0.80}$, $p^{0.97}$,
and $p^{0.90}$. Their margins below a full power of $p$ are $0.20$,
$0.03$, and $0.10$. For large $p$, the term $p^{0.97}$ grows fastest
and controls the exponent of the total.

Make the first operation much cheaper, changing $p^{0.80}$ to $p^{0.60}$.
Its margin improves to $0.40$, yet $p^{0.97}$ still controls the total.
Alternatively, improve the second operation to $p^{0.92}$. Now the minimum
margin is $0.08$, so the combined exponent can improve as well, subject
to its remaining overhead and strict-slack requirements.

\explfigure{22-bottleneck}{A local gain need not improve the overall exponent.}{
Each row is a separate toy profile, not a historical result. Larger margins
are better. The printed minimum stays at $0.03$ when an already easy term
improves, and rises to $0.08$ when the limiting term improves.}

The source's final accounting has seven dominant terms instead of three.
The figure in the next technical explanation computes their actual original
margins. Its logarithmic scale makes extremely small positive values visible.
That plot reverses the visual convention: farther right means a smaller
margin because the axis is minus the logarithm of the margin. Always read
the axis before deciding which bar is better.

\paragraph{Why a little extra room is needed.}
A cost may also carry factors such as $(\log p)^2$. Such a factor grows
more slowly than any fixed positive power $p^\rho$, but it is not a
constant. Choosing a final saving strictly below the available margin leaves
room to absorb these factors for sufficiently large $p$. Equality at the
boundary would not justify the same argument.

\takeaway{Track the full cost profile. An improvement is useful at the
headline level only when all the limiting terms and interfaces permit it.}
